\documentclass[11pt,a4paper]{article}
\usepackage{curso}
\begin{document}
\portada{Lección}{6.5}{Modelos de tiempo}

\noindent\textbf{Unidad 6.}\quad\textbf{Objetivos:} construir un modelo de tiempo de extremo a extremo (carga DAC, dinámica analógica, lectura ADC); entender
por qué la compilación $O(d^2)$ suele dominar cuando el núcleo analógico es rápido; amortizar la carga.
\textbf{Prerrequisitos:} 5.8, 6.1, 6.4.

\section{Tres etapas}
(1) \textbf{Compilación}: los $d^2$ elementos de $A$ pasan por convertidores digital-analógico (DAC). (2) \textbf{Dinámica}: equilibrar y promediar,
$t_0+\tau\propto d\,\tau_r$ para el sistema lineal. (3) \textbf{Lectura}: los $d$ valores de $x$ pasan por convertidores analógico-digital (ADC). Los métodos
digitales no pagan (1): la matriz ya está en su memoria.

\section{Constante \texorpdfstring{$RC$}{RC}}
En un circuito, $\tau_r=RC$. Con $1\ \mu$s el núcleo es lento frente a la carga; con $10$ ns (típico de un circuito integrado pequeño \evE) es cien veces
más rápido. Cuanto más rápido el núcleo analógico, más pesa la parte digital.

\section{Modelo explícito \evE}
\[T(d)=\frac{d^2}{B_{\text{DAC}}}+c\,d\,\tau_r+\frac{d}{B_{\text{ADC}}},\qquad B=10^9\ \text{valores/s},\ c=20\]
($c$ de la escala de la simulación de la lección 6.1). Para $d=1000$:
\begin{center}\begin{tabular}{lccc}\toprule
& carga & dinámica & fracción de carga\\\midrule
$RC=1\ \mu$s & 1,0 ms & 20 ms & 5\,\%\\
$RC=10$ ns & 1,0 ms & 0,2 ms & 83\,\%\\\bottomrule\end{tabular}\end{center}
Cruce carga = dinámica: $d^\ast=c\,\tau_rB$, es decir $2\times10^4$ ($1\ \mu$s) o $200$ ($10$ ns).

\section{Amortización}
Con $K$ lados derechos y la misma $A$: $T_K/K=d^2/(KB)+c\,d\,\tau_r+d/B$. Para $d=1000$, $RC=10$ ns: $K=1$: 1,20 ms (carga 83\,\%); $K=10$: 0,30 ms (33\,\%);
$K=100$: 0,21 ms (5\,\%). En un método de Newton, donde $A$ cambia en cada paso, la compilación se paga siempre.

\section{Lo que encontró el artículo}
Aifer et al.\ incluyen todas las operaciones digitales, de conversión y analógicas y comparan con gradiente conjugado y Cholesky en GPU: con $d=100$ el método
termodinámico pierde; con $d=1000$ y $5000$ es competitivo, sobre todo a tiempos cortos. Su modelo aún no incluye la imprecisión analógica.

\begin{aplicacion}
Reutilizar $A$ (procesos gaussianos, precondicionadores fijos); actualizaciones de bajo rango en vez de recargar $d^2$ elementos; bancos paralelos de DAC y memoria de
acoplamientos en el chip. Laboratorio L9: añadir un modelo de tiempo y encontrar $d^\ast$.
\end{aplicacion}

\begin{error}
\begin{itemize}[nosep]
\item Reportar solo el tiempo analógico. \quad Suponer que un núcleo más rápido siempre ayuda.
\item Olvidar que la línea base digital no paga transferencia. \quad Recargar $A$ cuando solo cambió $b$.
\end{itemize}
\end{error}

\begin{thebibliography}{9}\small
\bibitem[30]{tla} M. Aifer et al., Thermodynamic linear algebra, \emph{npj Unconv. Comput.} 1, 13 (2024).
\bibitem[25]{mel} D. Melanson et al., Thermodynamic computing system for AI applications, \emph{Nat. Commun.} 16, 3757 (2025).
\bibitem[29]{asic} M. Aifer et al., Solving the compute crisis with physics-based ASICs, 2025.
\end{thebibliography}
\end{document}
